\section{Problem Formulation and Boundary Conditions}
\label{sec:problem_formulation_mode1}

The baseline verification problem represents a square specimen of dimensions $L \times W = 1.0\,\text{mm} \times 1.0\,\text{mm}$ subject to tensile Mode-I loading, adopting the benchmark geometry of Pandey and Kumar~\citep{pandey2025}. An initial horizontal sharp crack of length $a_0 = 0.5\,\text{mm}$ is embedded along the centerline ($y = 0.5\,\text{mm}$, $0 \le x \le 0.5\,\text{mm}$).

\paragraph{Kinematic Boundary Conditions:}
\begin{itemize}
    \item \textbf{Bottom Boundary ($y = 0\,\text{mm}$)}: Roller constraint preventing vertical displacement, $u_y(x, 0) = 0$.
    \item \textbf{Bottom-Left Corner Pin ($(x, y) = (0, 0)$)}: Rigid-body horizontal constraint, $u_x(0, 0) = 0$.
    \item \textbf{Top Boundary ($y = 1.0\,\text{mm}$)}: Homogeneous vertical tensile displacement $u_y(x, 1) = u(t)$ prescribed via a reference point (RP) kinematically coupled to the top edge nodes, with free lateral contraction ($u_x$ unconstrained).
    \item \textbf{Lateral Boundaries ($x = 0\,\text{mm}, x = 1.0\,\text{mm}$)}: Traction-free boundaries, $\boldsymbol{\sigma} \cdot \mathbf{n} = \mathbf{0}$.
\end{itemize}

\paragraph{Constitutive Material and Phase-Field Properties:}
The material is modeled under plane-strain linear elasticity with AT2 phase-field regularized fracture:
\begin{itemize}
    \item Young's Modulus: $E = 210\,\text{GPa} \equiv 210.0\,\text{kN/mm}^2$
    \item Poisson's Ratio: $\nu = 0.3$
    \item Critical Fracture Energy Release Rate: $G_c = 2.7\times 10^{-3}\,\text{kN/mm} \equiv 2.7\,\text{N/mm} = 2700\,\text{J/m}^2$
    \item Phase-Field Regularization Length Scale: $l_0 = 0.0075\,\text{mm} \equiv 7.5\,\mu\text{m}$
    \item Residual Degradation Parameter: $k_{\text{res}} = 1.0 \times 10^{-7}$
\end{itemize}

\section{Crack Geometry Representation: Sharp Seam vs. Blunt Notch}
\label{sec:seam_vs_notch_mechanics}

A fundamental mechanical distinction in fracture benchmarks is the representation of initial crack geometry. While literature implementations frequently approximate cracks as finite-width blunt notches ($r_0 > 0$) to avoid singular meshing, the benchmark formulation requires an ideal sharp crack ($a_0 = 0.5\,\text{mm}$) modeled as a zero-gap duplicate-node seam.

\paragraph{Global Structural Stiffness Definition:}
The initial structural stiffness $K_0$ is defined as the slope of the linear-elastic reaction force ($F$) versus prescribed top-edge displacement ($u$) curve before damage initiation:
\begin{equation}
K_0 = \left. \frac{\mathrm{d}F}{\mathrm{d}u} \right|_{\text{initial elastic range}} \approx \frac{\Delta F}{\Delta u}.
\end{equation}
While Young's modulus $E$ is a local constitutive tensor property ($210\,\text{kN/mm}^2$), $K_0$ represents the integrated load-deformation compliance of the entire boundary value problem, depending directly on specimen geometry, boundary constraints, and notch kinematics.

Early geometry-sanity-check models (Figure~\ref{fig:notch_stiffness_thesis}) compared an uncracked solid ($K_0 \approx 138.15\,\text{kN/mm}$), a blunt notch ($K_0 \approx 75.47\,\text{kN/mm}$), and an early sharp seam diagnostic ($K_0 \approx 141.82\,\text{kN/mm}$), demonstrating that finite-width blunt notches introduce artificial compliance by removing elastic material. These values represent early sanity-check configurations under preliminary boundary approximations; the final governed canonical $S_1$ benchmark stiffness established below is $K_0 = 137.945520\,\text{kN/mm}$.

\begin{figure}[htbp]
\centering
\includegraphics[width=0.75\textwidth]{fig_notch_stiffness.pdf}
\caption[Early geometry-sanity-check comparison of structural stiffness $K_0$.]{Early geometry-sanity-check comparison of structural stiffness $K_0$: uncracked solid ($\approx 138.15\,\text{kN/mm}$), blunt notch introducing artificial compliance ($\approx 75.47\,\text{kN/mm}$), and early sharp seam diagnostic ($\approx 141.82\,\text{kN/mm}$). The final governed canonical $S_1$ sharp-seam reference stiffness is $K_0 = 137.945520\,\text{kN/mm}$ (Section~\ref{sec:fixed_mesh_reference_anchor}).}
\label{fig:notch_stiffness_thesis}
\end{figure}

In contrast to blunt notches, the zero-gap sharp seam models the crack as duplicate coincident node lines with identical initial spatial coordinates ($y = 0.5^+\,\text{mm}$ and $y = 0.5^-\,\text{mm}$ for $0 \le x \le 0.5\,\text{mm}$). This preserves the full continuous elastic volume while permitting discontinuous separation.

\section{Multi-Layer UEL/UMAT Layered Architecture in Abaqus}
\label{sec:uel_architecture_abaqus}

To implement the staggered phase-field formulation in Abaqus/Standard, a three-layer co-located element architecture is constructed:
\begin{itemize}
    \item \textbf{Layer 1 (Phase-Field Damage Layer)}: Co-located 4-node user elements (\texttt{UEL}, user element type \texttt{U1}, \texttt{JTYPE = 1}) solving for the scalar damage degree of freedom $d \in [0, 1]$ mapped to Abaqus \texttt{DOF 3} (variable \texttt{U3}) governed by the regularized AT2 Helmholtz-type PDE.
    \item \textbf{Layer 2 (Mechanical Displacement Layer)}: Co-located 4-node quadrilateral user elements (\texttt{UEL}, user element type \texttt{U2}, \texttt{JTYPE = 2}) solving for displacement degrees of freedom $(u_x, u_y)$ (\texttt{DOF 1} and \texttt{DOF 2}) using the degraded isotropic elastic tangent operator $g(\bar{d})\mathbb{C}_0$.
    \item \textbf{Layer 3 (Companion Visualization UMAT Layer)}: Standard continuum plane-strain elements (\texttt{CPE4}) sharing identical connectivity and nodal coordinates, driven by a dummy user material (\texttt{UMAT}) that maps internal state variables (SDVs) into standard Abaqus field outputs (\texttt{SDV1} = Phase Field $d$, \texttt{SDV2} = History $\mathcal{H}$, \texttt{SDV15} = degraded stiffness factor $g(d)$, \texttt{SDV17} = whole-underlying-element fracture surface energy $E_{\mathrm{frac}}^{(e)}$ in native $\mathrm{kN\cdot mm} = \mathrm{J}$, \texttt{SDV18} = whole-underlying-element elastic strain energy $E_{\mathrm{elas}}^{(e)}$ in native $\mathrm{kN\cdot mm} = \mathrm{J}$, with reporting in $\mathrm{mJ}$ performed during post-processing via a single $1000\times$ multiplication factor; \texttt{SDV19}/\texttt{SDV20} = area-normalized source quotients $\bar{\psi}_f$ and $\bar{\psi}_e$ in $\mathrm{kN/mm} = \mathrm{J/mm^2}$, interpretable as volumetric energy densities only under the explicit $1\text{-}\mathrm{mm}$ implicit unit-thickness convention) for visualization in Abaqus/Viewer.
\end{itemize}

All three layers communicate state variables across solver iterations via global Fortran common blocks in user subroutine \nolinkurl{f42_mixed_uel.for} (\nolinkurl{SHA256: CE8D5EDCD2911DCB018BB15275271F874E7EA62B8FB48CF4A8297469A83ACDD6}).

\subsection{Discrete Weak-Form Energy Derivation and Deduplication Proof}
\label{sec:uel_energy_derivation}

In user subroutine \nolinkurl{f42_mixed_uel.for}, stored elastic strain energy $E_{\mathrm{elas}}$ and regularized fracture surface energy $E_{\mathrm{frac}}$ are evaluated over the $N_{\mathrm{mesh}}$ underlying finite elements via standard $2\times 2$ Gauss quadrature:
\begin{align}
  E_{\mathrm{elas}} &= \sum_{e=1}^{N_{\mathrm{mesh}}} \int_{\Omega_e} \frac{1}{2} g(\bar{d}_e)\,\boldsymbol{\varepsilon} : \mathbb{C}_0 : \boldsymbol{\varepsilon} \,\mathrm{d}\Omega = \sum_{e=1}^{N_{\mathrm{mesh}}} \sum_{k=1}^{4} \frac{1}{2} g(\bar{d}_e)\,\boldsymbol{\varepsilon}(\boldsymbol{\xi}_k) : \mathbb{C}_0 : \boldsymbol{\varepsilon}(\boldsymbol{\xi}_k) \, \det(\mathbf{J}_k) W_k \, B, \label{eq:elas_weak_energy} \\
  E_{\mathrm{frac}} &= \sum_{e=1}^{N_{\mathrm{mesh}}} \int_{\Omega_e} G_c \left[ \frac{\bar{d}_e^2}{2 l_0} + \frac{l_0}{2} |\nabla d|^2 \right] \mathrm{d}\Omega = \sum_{e=1}^{N_{\mathrm{mesh}}} \sum_{k=1}^{4} G_c \left[ \frac{\bar{d}_e^2}{2 l_0} + \frac{l_0}{2} |\nabla d(\boldsymbol{\xi}_k)|^2 \right] \det(\mathbf{J}_k) W_k \, B, \label{eq:frac_weak_energy}
\end{align}
where $\bar{d}_e = \frac{1}{4}\sum_{a=1}^4 d_a$ is the element-averaged damage, $g(\bar{d}_e) = (1-\bar{d}_e)^2 + k_{\mathrm{res}}$, and $B = 1.0\,\mathrm{mm}$ is the unit plane-strain thickness.

External mechanical boundary work $W_{\mathrm{trap}}$ is evaluated from the reaction force $F = RF_2$ at the reference point:
\begin{equation}
  W_{\mathrm{trap},n} = \sum_{i=1}^{n} \frac{F_i + F_{i-1}}{2}(u_i - u_{i-1}), \quad \text{with } u_0 = 0, F_0 = 0.
\end{equation}

\paragraph{Analytical Deduplication and Companion Neutrality:}
A critical consideration in co-located multi-layer element architectures is the prevention of artificial double accumulation of internal energies. In the present formulation:
\begin{enumerate}
    \item \textbf{Layer 1 (Phase UEL)}: Evaluates only $E_{\mathrm{frac}}^{(e)}$ and assigns \texttt{ENERGY(7) = E\_FRAC\_ELEM}. It possesses zero mechanical displacement DOFs, zero stress ($\boldsymbol{\sigma} = \mathbf{0}$), and computes zero elastic strain energy ($E_{\mathrm{elas}} = 0$).
    \item \textbf{Layer 2 (Mechanical UEL)}: Evaluates only $E_{\mathrm{elas}}^{(e)}$ and assigns \texttt{ENERGY(2) = E\_ELAS\_ELEM}, which Abaqus sums natively into \texttt{ALLSE}. It contains zero damage DOFs and performs zero fracture functional computation.
    \item \textbf{Layer 3 (Companion UMAT)}: Is assigned an infinitesimal dummy elasticity tensor $\mathbf{D}_{\mathrm{comp}} = 10^{-11}\,\mathbf{I}\,\mathrm{GPa}$, producing identically vanishing physical stresses ($\boldsymbol{\sigma}_{\mathrm{comp}} \approx \mathbf{0}$) and zero internal energy contributions ($\text{SSE} = \text{SPD} = \text{SCD} = 0$, contributing $\sim 10^{-11}\,\mathrm{J} \approx 0$).
\end{enumerate}
Consequently, each energetic physical domain is evaluated on exactly one layer, guaranteeing no counted physical-energy duplication.

\paragraph{History-Field Non-Potentiality and Bookkeeping Residual:}
Because damage irreversibility is enforced through the non-decreasing historical maximum $\mathcal{H}(\mathbf{x}, t) = \max_{\tau \le t} \psi_0^+(\boldsymbol{\varepsilon}(\tau))$, the governing system does not represent the stationary point of any single instantaneous potential $\Pi(\mathbf{u}, d)$. The discrete bookkeeping residual $\Delta_{\mathrm{book}} = W_{\mathrm{trap}} - (E_{\mathrm{elas}} + E_{\mathrm{frac}})$ is monitored throughout deformation. For the canonical $S_1$ reference baseline (Job \texttt{1409734.mmaster02}), $\Delta_{\mathrm{book}} = \mathbf{+0.017948\,\mathrm{mJ}}$ ($\varepsilon_{\mathrm{book}} = \mathbf{0.7607\%}$) at complete separation ($u = 0.010\,\text{mm}$), demonstrating near-complete energy conservation within the unclosed rate-dissipation boundary.

\subsection{Mechanical Non-Invasiveness \& Full-History Parity Qualification}
\label{sec:mechanical_parity_qualification}

To prove that energy instrumentation in user subroutine \nolinkurl{f42_mixed_uel.for} does not alter the underlying mechanics, a full-history point-by-point parity comparison was conducted between the pre-instrumentation fixed reference solve (\nolinkurl{01_standard_pfm_reference}, Fortran SHA-256 \nolinkurl{ED1586D6427A4B1A01D99F7E219891EC7BE9FE911E066D9360724942E7D27720}) and the energy-instrumented fixed reference solve (\nolinkurl{16_energy_qualification_reference_15k}, Job \texttt{1409734.mmaster02}, Fortran SHA-256 \nolinkurl{CE8D5EDCD2911DCB018BB15275271F874E7EA62B8FB48CF4A8297469A83ACDD6}) across all $7{,}000$ increments:
\begin{itemize}
    \item \textbf{Initial Stiffness}: $K_0 = 137.945519645\,\text{kN/mm}$ for both models ($\Delta K_0 = 0.000000\,\text{kN/mm}$, $0.000000\%$ difference, bitwise match to 12 decimal places).
    \item \textbf{Peak Load}: $F_{\max} = 0.75777849\,\text{kN}$ at $u = 0.005857\,\text{mm}$ ($\Delta F_{\max} = 0.000000\,\text{kN}$, $\Delta u = 0.000000\,\text{mm}$).
    \item \textbf{Solver Convergence}: Completed increments ($7{,}000/7{,}000$), Newton equilibrium iterations ($21{,}120/21{,}120$), cutbacks ($0/0$), and severe discontinuities ($0/0$) are 100\% identical.
    \item \textbf{External Boundary Work}: $W_{\mathrm{trap}} = 2.359328928\,\text{mJ}$ vs $2.359328928\,\text{mJ}$ ($\Delta W = 7.12\times 10^{-11}\,\text{mJ}$, $3.02\times 10^{-9}\%$ difference, representing pure numerical roundoff).
\end{itemize}
This establishes the formal promoted status: \textbf{\texttt{UEL\_ENERGY\_OUTPUT\_QUALIFIED\_MECHANICALLY\_NONINVASIVE}}.

\section{Canonical S1 Reference Baseline Provenance}
\label{sec:fixed_mesh_reference_anchor}

The governed canonical reference baseline for all subsequent evaluations is Job \texttt{1406015.mmaster02} / Job \texttt{1409734.mmaster02} ($S_1$, $15{,}192$ underlying finite elements, $15{,}522$ total \texttt{*NODE} entries ($15{,}521$ mesh nodes + RP 999999)), generated from the canonical input deck (\nolinkurl{SHA256: bc000e0c5791150a711cb7af447a0bb34d1b9d9af5c6ce9d096c98749e8edf61}).

\begin{table}[htbp]
\centering
\caption{Quantitative response metrics of the Mode-I canonical S1 baseline.}
\label{tab:reference_anchor_metrics}
\small
\begin{tabular}{lrr}
\toprule
\textbf{Metric / Parameter} & \textbf{Canonical S1 Baseline (Job 1406015.mmaster02)} & \textbf{Literature \citep{pandey2025}} \\
\midrule
Initial Structural Stiffness $K_0$ & $137.945520\,\text{kN/mm}$ & Project-derived baseline \\
Linear Intercept $b$ ($u \le 1.0\,\mu\text{m}$) & $4.472368 \times 10^{-5}\,\text{kN}$ & --- \\
Coefficient of Determination $R^2$ & $0.99999960$ & --- \\
Peak Reaction Force $F_{\max}$ & $0.757778\,\text{kN}$ & $\approx 0.758\,\text{kN}$ \\
Displacement at Peak $u(F_{\max})$ & $0.005857\,\text{mm}$ ($5.857\,\mu\text{m}$) & $\approx 0.005860\,\text{mm}$ \\
Initial Crack Seam Length $a_0$ & $0.5000\,\text{mm}$ & $0.50\,\text{mm}$ \\
Crack Extension Symmetry & $y = 0.5000\,\text{mm}$ ($\Delta y \le 0.0031\,\text{mm}$) & Pure Mode-I planar \\
\bottomrule
\end{tabular}
\end{table}

As summarized in Table~\ref{tab:reference_anchor_metrics}, unconstrained ordinary least-squares regression on the audited data source \nolinkurl{S1_h0030_15k_MECHANICAL_FU_AUDITED.csv} (\nolinkurl{SHA256: E3100B7429E14B8BCEAB9952DDF8B0419CA9F71DF058F5916668922185AA25E2}) over the initial $N = 400$ active increments ($u \le 1.0\,\mu\text{m}$) yields $K_0 = 137.945520\,\text{kN/mm}$ with $R^2 = 0.99999960$. The peak reaction force $F_{\max} = 0.757778\,\text{kN}$ occurs at $u = 0.005857\,\text{mm}$, matching the digitized literature response of Pandey and Kumar~\cite{pandey2025} ($F_{\max} \approx 0.758\,\text{kN}, u \approx 0.005860\,\text{mm}$). Historical Job \texttt{1398090.mmaster02} is preserved as reconstructed baseline evidence.
